% ===== 3. Sourcing and curation of the ConfabQA question set =====
\section{3. Sourcing and curation of the ConfabQA question
set}\label{sourcing-and-curation-of-the-confabqa-question-set}

ConfabQA comprises 784 items in a \(4 \times 3\) design: 4 domains
(science, history, culture, cinema) crossed with 3 categories
(well-known pre-cutoff, obscure pre-cutoff, post-cutoff). Across domains
the categories hold 143 well-known, 153 obscure, and 488 post-cutoff
items (per-cell counts in Table \ref{tbl:cells}, Appendix A). I generate
items from per-domain templated source files. Each carries a provenance
URL for its gold answer and a validation-status field. The 384 items marked verified were marked
from source-file provenance only, with no independent check of the gold answer. The other 400, all
post-cutoff items from the dense post-cutoff source files, remain unverified. A wrong gold
answer can move an item between correct and wrong but not into or out of refusal, and the
web-verified regrade of all 784 labels (Section 4) gives an independent check. The paper analyzes the third of three
validation iterations, \textbf{ConfabQA} (\(n=784\)). The full design is
in \texttt{data/QUESTIONS\_v1\_CARD.md}, and pipeline reproducibility and
a sample source entry are in Appendix B.2 and B.5. A fifth domain
(sports) produced \(0/26\) correct, leaving no positive examples for the
within-pre-cutoff probe, so I exclude it (Appendix A.1).

The third (\texttt{obscure}) category populates the cell ``answer-was-in-training-data but
model-is-expected-to-fail,'' which the standard two-cell design leaves empty. A typical obscure item asks who won the Best Actor
Academy Award at the 1968 ceremony (gold: Rod Steiger, on the public
record). Qwen3-1.7B confidently answers with the wrong actor.
